\documentclass{exam-zh}

\examsetup{
  question/show-answer = false,
  fillin = {
    no-answer-type = none,
    width = auto
  }
}

\begin{document}

\section*{学生版：按答案自然宽度留空}

\begin{question}
  王湾《次北固山下》的名句“\fillin[海日生残夜]，\fillin[江春入旧年]”，
  描写时序交替中的景物，暗示着时光流逝，蕴含着自然理趣。
\end{question}

\begin{question}
  小慧为朋友家的农家乐餐厅写宣传横幅，直接使用了陆游《游山西村》里的
  “\fillin[山重水复疑无路]，\fillin[柳暗花明又一村]”两句诗。
\end{question}

\begin{question}
  行至群山深处，见到一挂瀑布飞泻而下，水石激荡，轰鸣作响，于老师回头说：
  “这不就是古诗中写的‘\fillin[飞流直下三千尺]，\fillin[疑是银河落九天]’嘛！”
\end{question}

\section*{教师版：相同输入}

\examsetup{question/show-answer = true}

\begin{question}
  王湾《次北固山下》的名句“\fillin[海日生残夜]，\fillin[江春入旧年]”，
  描写时序交替中的景物，暗示着时光流逝，蕴含着自然理趣。
\end{question}

\begin{question}
  小慧为朋友家的农家乐餐厅写宣传横幅，直接使用了陆游《游山西村》里的
  “\fillin[山重水复疑无路]，\fillin[柳暗花明又一村]”两句诗。
\end{question}

\begin{question}
  行至群山深处，见到一挂瀑布飞泻而下，水石激荡，轰鸣作响，于老师回头说：
  “这不就是古诗中写的‘\fillin[飞流直下三千尺]，\fillin[疑是银河落九天]’嘛！”
\end{question}

\end{document}
